\documentclass[12pt]{report}

% ============================================================
% Georgia State University Ph.D. Dissertation Proposal Template
% Department of Computer Science
% ============================================================

% ---------- Page setup ----------
\usepackage[letterpaper,margin=1in]{geometry}
\usepackage{setspace}
\doublespacing

% ---------- Fonts / text ----------
\usepackage[T1]{fontenc}
\usepackage[utf8]{inputenc}
\usepackage{lmodern}
\usepackage{microtype}

% ---------- Mathematics ----------
\usepackage{amsmath,amssymb,amsfonts,mathtools}
\usepackage{bm}

% ---------- Figures / tables ----------
\usepackage{graphicx}
\usepackage{booktabs}
\usepackage{multirow}
\usepackage{array}
\usepackage{caption}
\usepackage{subcaption}

% ---------- Lists / algorithms ----------
\usepackage{enumitem}
\usepackage[ruled,vlined]{algorithm2e}

% ---------- References / links ----------
\usepackage[numbers,sort&compress]{natbib}
\usepackage[hidelinks]{hyperref}
\usepackage[nameinlink,noabbrev]{cleveref}

% ---------- Misc ----------
\usepackage{xcolor}
\usepackage{url}
\usepackage{fancyhdr}

% ---------- Header / footer ----------
\pagestyle{fancy}
\fancyhf{}
\fancyhead[R]{\thepage}
\renewcommand{\headrulewidth}{0pt}

% ---------- Chapter formatting ----------
% Keep standard report chapter formatting for maximum portability.

% ---------- User information ----------
\newcommand{\ProposalTitle}{
Data intensive Learning and Ordinal Uncertainty Quantification
}

\newcommand{\StudentName}{Jinsu Hong}
\newcommand{\ProposalYear}{2026}

\newcommand{\AdvisorName}{Berkay Aydin, Ph.D.}

% Replace / add committee members as needed.
\newcommand{\CommitteeMemberA}{Jingyu Liu, Ph.D.}
\newcommand{\CommitteeMemberB}{Yue Wang, Ph.D.}
\newcommand{\CommitteeMemberC}{Talwinder Singh, Ph.D.}
\newcommand{\CommitteeMemberD}{Sujit Roy, Ph.D.}

\newcommand{\DepartmentChair}{Armin Mikler, Ph.D.}

\begin{document}

% ============================================================
% TITLE PAGE
% ============================================================

\begin{titlepage}
    \singlespacing
    \thispagestyle{empty}
    \enlargethispage{0.3in}
    
    \centering

    % ---------------------------------------------------------
    % Title
    % ---------------------------------------------------------
    \vspace*{0.10in}

    {\Large \ProposalTitle \par}

    \vspace{0.25in}

    A Dissertation Proposal

    \vspace{0.32in}

    Presented in Partial Fulfillment of Requirements for the Degree of

    \vspace{0.12in}

    Doctor of Philosophy

    \vspace{0.28in}

    in the

    \vspace{0.28in}

    College of Arts and Sciences

    \vspace{0.28in}

    Georgia State University

    \vspace{0.28in}

    \ProposalYear

    \vspace{0.28in}

    by

    \vspace{0.22in}

    \StudentName


    % ---------------------------------------------------------
    % Committee
    % ---------------------------------------------------------
    \vspace{0.35in}

    Committee:

    \vspace{0.18in}

    % Advisor / Chair
    \rule{4.5in}{0.4pt}\\[-0.02in]
    \AdvisorName, Chair

    \vspace{0.13in}

    % Member 1
    \rule{4.5in}{0.4pt}\\[-0.02in]
    \CommitteeMemberA, Member

    \vspace{0.13in}

    % Member 2
    \rule{4.5in}{0.4pt}\\[-0.02in]
    \CommitteeMemberB, Member

    \vspace{0.13in}

    % Member 3
    \rule{4.5in}{0.4pt}\\[-0.02in]
    \CommitteeMemberC, Member

    \vspace{0.13in}

    % Member 4
    \rule{4.5in}{0.4pt}\\[-0.02in]
    \CommitteeMemberD, Member


    % ---------------------------------------------------------
    % Date
    % ---------------------------------------------------------
    \vspace{0.18in}

    \rule{4.5in}{0.4pt}\\[-0.02in]
    Date


    % ---------------------------------------------------------
    % Department Chair
    % ---------------------------------------------------------
    \vspace{0.18in}

    \rule{4.5in}{0.4pt}\\[-0.02in]
    \DepartmentChair\\[-0.02in]
    Department Chair

\end{titlepage}

% ============================================================
% FRONT MATTER
% ============================================================
\pagenumbering{roman}

\chapter*{Abstract}
\addcontentsline{toc}{chapter}{Abstract}

% Keep the proposal abstract concise: motivation, dissertation theme,
% three research components, main methods, preliminary findings,
% and expected contributions.
%
% Example outline:
% 1. Scientific / methodological motivation
% 2. Solar-flare forecasting problem
% 3. Target design and multi-channel prediction
% 4. Ordinal uncertainty quantification
% 5. Expected dissertation contributions

\clearpage

\tableofcontents
\clearpage

\listoffigures
\addcontentsline{toc}{chapter}{List of Figures}
\clearpage

\listoftables
\addcontentsline{toc}{chapter}{List of Tables}
\clearpage

% ============================================================
% MAIN MATTER
% ============================================================
\pagenumbering{arabic}

% ============================================================
\chapter{Introduction}
\label{chap:introduction}

This chapter introduces the dissertation's three main research directions: prediction target design, learning from multi-channel scientific data, and uncertainty quantification for ordinal prediction. It presents the motivation, research problems, and research questions, followed by an overview of the dissertation and its expected contributions.

\section{Motivation}
Space weather refers to changing conditions in the space environment, primarily driven by solar activity, including changes in Earth's magnetosphere, ionosphere, and thermosphere \cite{noaa_space_weather}. Among the various forms of solar activity, solar flares are particularly important because they release large amounts of electromagnetic radiation, including X-rays and extreme ultraviolet radiation \cite{noaa_solar_flares}. Solar flares and other solar eruptions, such as coronal mass ejections (CMEs), can have significant impacts on Earth and near-Earth space. These impacts include disruptions to radio communications, satellite navigation systems, and power grids. In addition, increased radiation levels during solar radiation storms can pose health risks to astronauts and people traveling at high altitudes \cite{noaa_space_weather_impacts}. Consequently, reliable forecasting of solar flares and related space weather events is important to mitigate their potential impacts.

Reliable space weather forecasting remains challenging for several reasons. Solar observations are collected by instruments with limited spatial and temporal resolution, and the resulting data may contain noise, missing observations, or measurement errors. Earth-facing observatories such as the Solar Dynamics Observatory (SDO) image the visible solar disk; activity on the far side of the Sun cannot be directly observed in these images. Severe flares are also rare relative to quiet or weakly active periods, leaving few examples of the events that matter most and creating imbalanced forecasting datasets \citep{barnes2016comparison,pelkumdonahue2024transformer}. In addition, the physical processes that trigger flares are not fully understood, which limits how completely observed precursors can be linked to future events \citep{pelkumdonahue2024transformer}. These constraints motivate careful choices of observations, forecast targets, and evaluation protocols; they do not imply that any one modeling approach resolves the underlying solar physics.

Many solar flare forecasting methods predict whether a flare exceeding a given intensity threshold will occur within a future time window. Such predictions provide a clear warning about potentially significant events, but they do not indicate how far the predicted flare intensity lies above or below the threshold. They also group events with different intensities into the same category. Previous work explored alternative labels based on maximum flare intensity, relative X-ray flux increases and cumulative flare activity and evaluated them in classification tasks \citep{hong2023innovative,hong2023beyond,hong2023enhancing}. This dissertation extends that work by studying continuous X-ray outcomes and examining how different target definitions affect predictability.

Forecasting performance also depends on the observations provided to a model. Solar images and magnetic-field measurements describe different features of the Sun. Solar flare forecasting models can use different combinations of observation channels \citep{cinto2020framework,bobra2015solarflare}. Their performance may also be influenced by experimental choices, such as the forecasting window, image resolution, and model capacity. This dissertation examines whether combining multiple observation channels provides additional useful information for solar flare forecasting.

Beyond solar flare forecasting, many prediction problems involve ordered categories, such as disease severity levels, age estimations, or solar flare classes. Reporting only a single predicted category does not communicate which other categories may also be plausible. Previous studies have explored prediction sets that contain multiple plausible categories while providing statistical coverage guarantees \citep{lu2022ordinalaps}. However, providing reliable predictions for rare and extreme categories remains challenging. This dissertation investigates uncertainty quantification methods that produce prediction sets with coverage guarantees while accounting for the natural order of categories, particularly for rare and extreme classes.

Together, these studies investigate how target definitions, input observations, and uncertainty quantification affect the performance and reliability of machine learning models. The dissertation addresses complementary challenges in developing reliable machine learning methods for scientific data and ordered prediction problems.

\section{Research Problem}

Reliable prediction remains difficult when forecasting targets, input observations, and uncertainty estimates are not adequately evaluated. In solar flare forecasting, it remains unclear how continuous target definitions affect predictability and whether additional observation channels consistently improve forecasting performance. Furthermore, conventional prediction methods may not effectively represent uncertainty for rare and extreme classes in ordinal classification problems. This dissertation addresses these challenges through three complementary research directions: continuous target formulation, multi-channel solar image analysis, and uncertainty quantification for ordinal prediction.

\section{Research Questions}

The proposed dissertation addresses the following research questions:

\begin{enumerate}[label=\textbf{RQ\arabic*:}]
    \item How does the definition of a continuous future solar X-ray target affect
    its predictability from multi-channel solar images?

    \item How do different combinations of solar observation channels affect forecasting performance, and does adding more channels provide useful information or introduce redundant or noisy information?

    \item How can prediction sets for ordinal classification account for class order and provide reliable uncertainty estimates for rare and extreme classes?
\end{enumerate}

\section{Dissertation Overview}

Chapter~\ref{chap:background} provides background on solar flare forecasting, scientific observations, target definitions, and uncertainty quantification for ordinal classification. Chapter~\ref{chap:project1} focuses on continuous X-ray target formulation and predictability. Chapter~\ref{chap:project2} investigates the use of different solar observation channels for flare forecasting. Chapter~\ref{chap:project3} presents uncertainty quantification methods for ordinal classification, with an emphasis on rare and extreme classes. Chapter~\ref{chap:discussion} discusses the findings across the three research directions. Chapter~\ref{chap:timeline} outlines the remaining work and dissertation timeline, and Chapter~\ref{chap:conclusion} concludes the proposal.


\section{Expected Contributions}

The expected contributions of this dissertation are:

\begin{enumerate}
    \item A systematic formulation and evaluation of continuous solar X-ray targets to investigate how different target definitions affect predictability and their relevance to operational forecasting.

    \item An empirical analysis of how different combinations of solar observation channels affect forecasting performance, including whether additional channels provide useful predictive information.

    \item An uncertainty quantification method for ordinal classification that accounts for the natural order of categories and provides reliable prediction sets with class-conditional coverage guarantees, particularly for rare and extreme classes.

    \item An empirical evaluation of prediction performance and uncertainty across data-intensive solar flare forecasting and ordinal classification tasks.
\end{enumerate}

% ============================================================
\chapter{Background and Related Work}
\label{chap:background}

This chapter reviews the background and related work relevant to the dissertation. It first introduces solar flare forecasting, including X-ray measurements, forecasting targets, and existing prediction methods. It then reviews solar imaging data and machine learning approaches that use different observation channels. Finally, it introduces uncertainty quantification for ordinal classification, with a focus on conformal prediction, ordinal prediction sets, and class-conditional coverage.


\section{Solar Flare Forecasting and Target Definitions}

\subsection{Solar Flare Classification and GOES X-ray Measurements}

The Geostationary Operational Environmental Satellites (GOES), operated by the National Oceanic and Atmospheric Administration (NOAA), provide X-ray flux measurements that are widely used to detect and classify solar flares. The GOES X-ray Sensor (XRS) measures solar X-ray flux, allowing NOAA to identify flare events based on changes in the observed flux. In the NOAA event detection procedure, the start of a flare is identified by a sustained increase in X-ray flux, while the peak is determined by the maximum measured flux. The end of the event is identified when the flux decreases to a specified level relative to the peak and the pre-flare background \cite{noaa_goes_xray_flux}.

NOAA also maintains historical records of solar flares detected using GOES observations, with X-ray flare records dating back to 1975. These records include flare classes, peak X-ray flux, and temporal information such as start, peak, and end times. Additional information, including flare locations and associated NOAA active region numbers, is available for some events \cite{Angryk2020}. Solar active regions are areas of concentrated magnetic activity where many solar flares originate.

Solar flares are classified based on their observed peak X-ray flux in the $0.1$--$0.8\,\mathrm{nm}$ wavelength band. The GOES classification system divides solar flares into five major classes: A, B, C, M, and X \cite{Fletcher2011}. These classes follow a logarithmic scale, with the minimum X-ray flux of each successive class being ten times higher than that of the previous class. X-class flares have peak fluxes of at least $10^{-4}\,\mathrm{W\,m^{-2}}$ at Earth. Each class is further divided according to a numerical multiplier, typically ranging from 1.0 to 9.9, except for X-class flares, which can exceed X9.9. For instance, a C5.7 flare corresponds to a peak X-ray flux of approximately $5.7 \times 10^{-6}\,\mathrm{W\,m^{-2}}$.

\subsection{Forecasting Tasks and Target Definitions}

When constructing datasets for data-driven solar flare forecasting, a fixed-length future interval, known as the prediction window, is commonly used to define the forecasting target. Within this window, the maximum flare intensity is often used to assign a categorical label, while a flare-quiet label may be assigned if no flare is detected. Binary classification tasks commonly use thresholds such as M1.0 or C5.0 to distinguish between flaring and non-flaring instances. The assigned label depends on whether the maximum flare intensity within the prediction window exceeds the selected threshold.

In addition to classification, solar flare forecasting can be formulated as a regression problem, where the target is a continuous measurement of future solar activity. Unlike categorical labels, continuous targets preserve differences in X-ray flux values. However, different continuous targets may represent different aspects of solar activity, such as maximum flux, accumulated flux, or changes relative to background conditions. Therefore, the choice of target determines which aspect of future solar activity a model is trained to predict.

\subsection{Related Work on Solar Flare Prediction Targets}

Many previous studies have formulated solar flare forecasting as a classification problem. These studies typically predict whether a flare exceeding a specified intensity threshold will occur within a future prediction window. Binary classification targets, such as the occurrence of M-class or stronger flares within 24 hours, have been widely used to evaluate flare forecasting methods \citep{barnes2016comparison,cinto2020framework, pandey2025ordinal,pandey2024embedding,ji2024towards,hong2023enhancing,hong2023innovative}.

Other studies have formulated solar flare forecasting as a regression problem to predict continuous measures of future solar activity. Muranushi et al. \citep{muranushi2015ufcorin} developed a regression-based forecasting system to predict the maximum GOES soft X-ray flux over the next 24 hours. Their target was defined as the maximum value of the continuous GOES X-ray flux measurements within the prediction window. Jiao et al. \citep{jiao2020solar} investigated the prediction of maximum flare intensity using GOES flare records and magnetic parameters from solar active regions. They used 24-hour prediction windows with different forecast horizons, including 0--24, 6--30, 12--36, and 24--48 hours. Although both studies predicted continuous flare intensity, their target definitions and prediction units differed.

Although previous studies have investigated continuous regression targets, the choice of target remains an important consideration. For example, maximum X-ray flux represents the strongest observed activity within a prediction window, while cumulative flux and changes relative to background conditions describe other characteristics of solar activity. Previous work also explored alternative flare labels based on relative X-ray increases and cumulative activity \citep{hong2023innovative,hong2023beyond,hong2023enhancing}. However, differences in target definitions and evaluation settings make it difficult to directly compare the predictability of different targets across studies. This dissertation investigates how target definitions affect predictability by evaluating multiple continuous X-ray targets under a common experimental setting.

\section{Solar Observations and Data-Driven Forecasting}

\subsection{Solar Imaging and Magnetic Field Observations}

In the field of solar flare forecasting, two main approaches are commonly used. The first one is referred to as active region-based models, which focus on data from active regions, represented as point-in-time vectors\cite{ahmadzadeh2019challenges}, time series \cite{ji2022solar}, or images\cite{abed2021automated}. In this approach, time series classification techniques are often employed, and flares are labeled based on their association with specific active regions \cite{ji2020all}. The second approach involves full-disk models, which provide predictions for the entire solar disk. In full-disk models, all flares are considered, regardless of their association with active regions \cite{pandey2021solar}. Hybrid approaches also exist in the literature \cite{pandey2022towards}. In these approaches, the primary objective is to predict the occurrence of significant flaring events within a defined prediction window, typically spanning 12, 24, or 48 hours.

In both active region-based \cite{ji2020all} or full-disk \cite{pandey2021solar} predictions, each instance is labeled with the maximum intensity flare in the prediction window, and previous studies in solar flare prediction primarily focus on labels derived from GOES classification, which can overlook pertinent information within the prediction windows. In \cite{Hong2023}, we presented a framework to create data-driven flare indices and introduced an extended dataset of flares with their associated data-driven labels. In this study, we employ these labels to predict flaring activity under both active region-based and full-disk flare prediction modes with different threshold settings. As part of our methodology, we integrate new machine learning-ready datasets and create new individual and cumulative labels associated with both active regions and full-disk data instances. We explore the performance and operational feasibility of these new labels with active region-based and full-disk predictions, considering a wide spectrum of threshold values and class weights. 

\subsection{Machine Learning for Solar Flare Forecasting}

Machine-learning (ML) flare forecasts learn a mapping from observations available at a forecast issue time to an event label or other target defined over a future window. As described above, studies vary in their prediction units, labels, and forecast windows; they also differ in how they represent the solar observations supplied to a model. Early ML systems commonly summarized active-region magnetograms with manually defined physical parameters before fitting a classifier. Bobra and Couvidat, for example, used an SVM with 25 parameters derived from SDO/HMI vector magnetic-field data for active-region flare prediction \citep{bobra2015solarflare}.

Deep-learning methods have used both image and time-series representations. Huang et al. trained a CNN directly on line-of-sight magnetograms from SOHO/MDI and SDO/HMI, allowing the model to learn image features without a separate hand-crafted feature-extraction stage \citep{huang2018flarelos}. By contrast, Liu et al. used an LSTM to model temporal sequences of SHARP magnetic parameters and flare-history features; their study was not a CNN operating on raw magnetogram images \citep{liu2019flarelstm}. Park et al. applied CNNs to full-disk magnetograms, using SOHO/MDI observations for earlier years and SDO/HMI observations for later years \citep{park2018fulldiskcnn}. These examples illustrate that architecture and input representation are separate design choices: CNNs can learn spatial patterns from image data, while recurrent models can represent temporal sequences of engineered features.

Deep Flare Net (DeFN) is an example of a deep neural network built on extracted active-region features. Its features combine magnetic-field information with ultraviolet, extreme-ultraviolet, and X-ray measurements, and the operational version uses near-real-time data products \citep{nishizuka2018defn,nishizuka2021defn}. Because these systems differ in whether they model image pixels, derived physical features, or feature sequences, their results should be interpreted together with the data source, forecast unit, target definition, and evaluation split. A score difference alone does not isolate the effect of the model architecture.

\subsection{Multi-channel Learning and Related Work}

Solar observations provide complementary measurements rather than interchangeable image channels. HMI measures photospheric magnetic fields, including line-of-sight and vector products; AIA records ultraviolet and extreme-ultraviolet images at multiple wavelengths; and GOES XRS measures disk-integrated soft X-ray flux. Combining these sources can expose a model to magnetic structure, atmospheric emission, and X-ray activity, but it also requires the inputs to be aligned in time and space and made available under the intended forecasting latency \citep{nasa_sdo_instruments}.

Jonas et al. developed image-derived features from SDO observations spanning the photosphere, chromosphere, transition region, and corona, and combined them with flare-history and physically motivated features for active-region prediction \citep{jonas2018flare}. Their results showed that the utility of additional inputs depended on the evaluation metric: photospheric vector-field data combined with flare history performed best for TSS, whereas the full collection of inputs helped for area under the precision-recall curve. DeFN provides another multi-source example. Its operational input includes HMI line-of-sight and vector magnetograms, AIA 131~\AA{} and 1600~\AA{} measurements, and GOES 1--8~\AA{} X-ray emission. The AIA 193~\AA{} and 304~\AA{} images are also shown for context in the operational system, so they should not be described as model inputs without specifying the version and feature set \citep{nishizuka2018defn,nishizuka2021defn}.

Recent heliophysics foundation models extend this direction beyond a single forecasting task. Surya is a spatiotemporal foundation model pretrained on multi-instrument SDO observations, including eight AIA channels and five HMI products; its reported downstream tasks include solar-flare forecasting \citep{roy2025surya}. It is therefore a multi-instrument model, not a single-channel AIA 171~\AA{} or HMI baseline. More generally, channel count alone does not specify a model's information content: the selected observables, their temporal context, the fusion procedure, and the prediction target all matter. Project~2 addresses these choices by comparing input channels alongside forecast windows, image resolution, and model capacity.


\section{Uncertainty Quantification for Ordinal Classification}

\subsection{Ordinal Classification and Prediction Uncertainty}

Ordinal classification predicts a label from an ordered set of categories, such as increasing levels of disease severity or solar flare class. The ordering carries information that a nominal classifier does not represent: predicting an adjacent class and predicting a distant class are different types of error. A single predicted label, however, does not show how many outcomes remain plausible. Set-valued prediction addresses this limitation by returning a subset of labels that is intended to contain the true outcome at a specified rate. For an input $x$, write this set as $C(x)\subseteq\mathcal{Y}$, where $\mathcal{Y}=\{0,1,\ldots,K-1\}$ is ordered from the lowest to the highest category.

The structure of a prediction set matters for ordinal outcomes. A set containing two separated classes while omitting the classes between them is difficult to interpret as a range of plausible severities. An ordinal prediction set is contiguous when it includes every label between its smallest and largest members. Contiguity therefore gives an interval-like summary over the ordered labels, while the set size communicates how many classes remain plausible. These properties describe the output structure; they do not by themselves establish that the set contains the true class at a desired frequency. A statistical coverage procedure is needed for that guarantee.

\subsection{Conformal Prediction}

Conformal prediction (CP) provides a general way to construct prediction sets with finite-sample coverage guarantees. In split CP, the predictive model and any data-dependent choices are first fixed using training data. A separate calibration sample of $n$ examples is then used to compute nonconformity scores $S_i=s(X_i,Y_i)$, which quantify how atypical each calibration example is under the fitted model. For a target error rate $\alpha$, let $S_{(j)}$ denote the $j$-th order statistic of these scores augmented with one $+\infty$ value. The split-conformal threshold is
\begin{equation}
  \widehat q_{1-\alpha}=S_{(\lceil(n+1)(1-\alpha)\rceil)}.
\end{equation}
For ordinal classification, a score $s(x,k)$ for candidate label $k$ yields the set
\begin{equation}
  C_{1-\alpha}(x)=\{k\in\mathcal{Y}:s(x,k)\leq\widehat q_{1-\alpha}\},
\end{equation}
where $\widehat q_{1-\alpha}$ is calibrated on held-out examples. Under exchangeability of the calibration examples and a future test example, this construction provides marginal coverage,
\begin{equation}
  \Pr\{Y_{n+1}\in C_{1-\alpha}(X_{n+1})\}\geq 1-\alpha,
\end{equation}
where the probability averages over the calibration and test data \citep{vovk1999cp,romano2020classification}. The guarantee is distribution-free in the sense that it does not require a parametric model for the data distribution; it still depends on the exchangeability assumption and on keeping the calibration data separate from model fitting.

The choice of nonconformity score affects the size and composition of the resulting set, even when the coverage guarantee holds. For example, the Least Ambiguous set-valued Classifier uses the predicted probability of the candidate class, while Adaptive Prediction Sets accumulate probability mass according to the model's class ranking \citep{sadinle2019least,angelopoulos2020uncertainty,romano2020classification}. These methods can be applied to a fitted classifier, but their standard forms treat labels as nominal. In addition, marginal coverage is averaged over the population: it does not imply coverage conditional on a particular input or class. This distinction matters when class frequencies are highly unequal.

\subsection{Ordinal Conformal Prediction}

Ordinal conformal prediction methods incorporate the ordering of labels into the construction of prediction sets. Ordinal Adaptive Prediction Sets (OAPS) use the order of the classes to form contiguous sets from classifier probabilities \citep{lu2022ordinalaps}. This produces outputs that respect the ordinal structure, while retaining a marginal coverage objective. Conformal Prediction Sets for Ordinal Classification (COPOC) instead uses an ordinal probability construction based on unimodality to form contiguous prediction sets \citep{dey2023copoc}. These methods address the structural mismatch between nominal prediction sets and ordered outcomes, but their marginal coverage target does not guarantee the same coverage separately for each class.

Conformalized Quantile Regression (CQR) provides a related way to use numeric information. It combines quantile regression with conformal calibration to construct prediction intervals for a continuous response \citep{romano2019cqr}. When ordinal categories correspond to ranges of a numeric target, an interval can be related to the ordered label space through the category boundaries. A continuous interval is not, by itself, a discrete ordinal prediction set, however, and standard CQR does not establish class-conditional coverage for those discrete labels. Thus, ordinal prediction raises two distinct design requirements: the output should reflect the order of the labels, and its coverage should be assessed at the population level or separately across classes according to the task.

\subsection{Class-Conditional Coverage}

Marginal coverage can hide differences in performance across classes. If common classes are well covered and rare classes are undercovered, the overall coverage may still meet the nominal target because the aggregate is weighted by class prevalence. Class-conditional coverage instead requires that the true label be included at the target rate within each class. This is a group-level guarantee conditional on the label; it is not a guarantee conditional on every feature vector or every combination of observing conditions.

Mondrian conformal prediction targets group-specific coverage by dividing calibration examples into predefined groups and calibrating a separate threshold within each group \citep{vovk2005algorithmic,vovk2012conditional}. For class-conditional coverage, the groups are defined by the true class labels. Under exchangeability within each class, calibration scores from that class support a finite-sample coverage guarantee for future examples of the same class. The guarantee depends on this within-class exchangeability assumption and does not automatically extend to distribution shifts between calibration and deployment data.

The amount of calibration data in each group affects the informativeness of the result. With few examples in a class, the available score quantiles are coarse; if the requested finite-sample rank exceeds the number of finite calibration scores, a conservative procedure may return an unbounded threshold and include that class for every input. This preserves coverage but can increase prediction-set size. For an ordinal task, class-conditional calibration and contiguous-set construction therefore address complementary needs: the former targets reliability for each class, while the latter ensures that the returned labels form an interpretable range. Our accepted OCQR work combines these objectives by using true-label, class-specific calibration and an ordinal post-processing step that only adds labels to the raw prediction set \citep{hong2026ocqr}. Its coverage claim is tied to the stated within-class exchangeability and target-label consistency assumptions.


\section{Chapter Summary}
























\section{Solar Flare Forecasting and Target Labels}

\subsection{Space Weather and Solar Flare Forecasting}

Space weather refers to changing conditions caused by solar activity and their effects in the space around Earth. The Sun produces a steady flow of charged particles called the solar wind. It also produces solar flares and large clouds of magnetized material called coronal mass ejections, and can accelerate energetic particles. These events
can disturb radio communication, satellite operations, navigation, and power systems \citep{noaa_space_weather,noaa_space_weather_impacts}. National forecasting services monitor solar activity and issue forecasts, watches, and warnings for potential effects at Earth \citep{noaa_space_weather}.

Solar flares are one part of space weather. A flare releases energy in the Sun's atmosphere and can produce a rapid increase in soft X-ray radiation. This radiation provides a measurable signal for flare forecasting \citep{noaa_solar_flares}. The present dissertation studies forecasts of flare-related X-ray outcomes and flare classes. Broader forecasts, such as the arrival of coronal mass ejections or the timing of geomagnetic storms, are outside its scope.

The Geostationary Operational Environmental Satellites (GOES) carry an X-ray Sensor (XRS) that records this radiation. Flare classes are based on peak X-ray flux and are ordered from weaker to stronger events \citep{noaa_solar_flares}. These measurements support forecasts about both whether an event will occur and how strong the measured flux will become.

Many forecasting studies turn future activity into a category. A binary task asks whether activity will cross a chosen threshold during a forecast window. A class-prediction task assigns the window to a flare class. Both tasks make the result easy to state, but the label does not preserve every detail of the future flux. For example, two windows on the same side of a threshold receive the same binary label even if their peak fluxes differ. Values close to a threshold can also receive different labels when they fall on opposite sides. This follows from the label definition; it does not show that a classifier's errors are concentrated near that boundary.

\subsection{Forecast Targets and Window-based Labels}

Forecasting studies define their targets over an explicit future window, but the value assigned to that window varies by task. Common formulations include whether the maximum flare intensity crosses a class threshold, the class or peak X-ray flux of the strongest event, and an aggregate of activity across the window. For example, a recent benchmark compares labels based on maximum and cumulative flare intensity over a 24-hour window \citep{roy2026suryabench}. Studies also differ in forecast horizon, latency, event magnitude threshold, and whether the forecast unit is the whole visible Sun or an individual active region; these choices affect the meaning and comparability of reported performance \citep{barnes2016comparison,
cinto2020framework}. This dissertation uses the term ``forecast target'' for these window-defined quantities. It distinguishes them from a solar-physics flare index unless a cited study explicitly defines such an index.

GOES X-ray Sensor (XRS) measurements provide a common basis for detecting and
classifying flares by peak soft X-ray flux. A detection or cataloging rule
identifies an event from measured flux, whereas a forecasting target uses
observations available before a future window to predict an event or a
window-level quantity. These are related but distinct tasks. The continuous
XRS-B outcomes in Project 1, including maximum flux, fluence, and background
deviations, are direct measurements aggregated over the forecast window;
they are not interchangeable with thresholded event labels or a separately
defined flare index.

Earlier work in this project developed flare metadata and alternative labels
based on relative X-ray increases and cumulative activity. These studies
evaluated the labels as classification targets \citep{hong2023innovative,
hong2023beyond,hong2023enhancing}. They also discussed limits of using only
the maximum flare class, since that label leaves out background flux and
activity accumulated across a forecast window \citep{hong2023beyond,
hong2023enhancing}. The metadata and labels were released as a separate
dataset \citep{hong2023data}. Those studies provide the background for the
present work, but they did not evaluate the ten continuous regression targets
defined in Chapter~\ref{chap:project1}.

Regression keeps the target as a measured value rather than converting it to
an event category. The result depends on which value is chosen. A forecast of
the maximum flux answers a different question from a forecast of total flux
or change from a background level. Regression does not remove forecast error;
it makes the size of the predicted quantity part of the task. Chapter~\ref{chap:project1}
defines and compares the continuous targets used in this dissertation.

\section{Solar Observations and Image-based Forecasting}

Solar forecasting studies use measurements from several instruments. The
Solar Dynamics Observatory (SDO) provides images from the Atmospheric Imaging
Assembly (AIA) and magnetic-field observations from the Helioseismic and
Magnetic Imager (HMI). The GOES XRS instrument measures soft X-ray flux rather
than producing solar images. These data describe different aspects of solar
activity, so a forecasting system must align observations by time and define
which channels it will use \citep{nasa_sdo_instruments}.

Image-based studies commonly use either full-disk observations or
active-region-based inputs. A full-disk image covers the visible solar disk,
while an active-region approach identifies and models a localized region;
these designs therefore differ in their prediction unit and input context
\citep{cinto2020framework}. Because these are remote measurements made from a
spacecraft at a finite distance from the Sun, the observed data are also
constrained by instrument resolution and cadence, viewing geometry, and data
availability. In particular, a single Earth-facing view does not directly
show the far-side hemisphere. SDO provides full-disk HMI coverage and
multi-wavelength AIA images at high cadence, but full-disk coverage should not
be interpreted as simultaneous observation of the entire solar sphere
\citep{nasa_sdo_instruments,nasa_sun_farside}.

Image-based models learn patterns from these observations. A convolutional
neural network applies filters across nearby image pixels and combines the
resulting features across layers. Researchers have used such models with
solar images and magnetic-field maps for flare forecasting
\citep{bobra2015solarflare,jonas2018flare}. A
model can also receive several image channels together. The selected channels,
image resolution, and forecast window can change what information is available
to the model. Published results can be hard to compare when studies differ in
these choices, as well as in their data partitions and evaluation procedures
\citep{barnes2016comparison,cinto2020framework}.

Transformer-based models have also been applied to solar flare forecasting.
Some studies use attention-based image encoders for magnetograms; others
combine image, video, or time-series models
\citep{pelkumdonahue2024transformer}. This work reviews
those studies as part of the field, but does not evaluate a vision-transformer
backbone. Project 2 instead studies the effect of input channels, forecast
windows, image resolution, and model capacity within its planned comparison.

Solar flares, coronal mass ejections (CMEs), and solar energetic particles
(SEPs) are related space-weather phenomena, but they are distinct prediction
targets. Prior work has separately studied whether a flare is associated with
a CME or SEP \citep{inceoglu2018cme_sep}. This dissertation focuses on
flare-related X-ray outcomes and flare classes; it does not claim to forecast
CME or SEP occurrence.

\section{Uncertainty for Ordered Flare Classes}

A point forecast gives one predicted value or class. A probability forecast
assigns probabilities to possible classes. A prediction set lists classes
that remain plausible for an observation. Sets can help when a single class
would hide uncertainty, but their usefulness depends on both whether they
include the true class and how many classes they contain
\citep{sadinle2019least,romano2020classification}.

Conformal prediction uses a separate calibration sample to set a threshold
for forming prediction sets. Under an exchangeability assumption, standard
split conformal methods provide a finite-sample guarantee for overall, or
marginal, coverage \citep{vovk1999cp,romano2020classification}. This
guarantee concerns the test population as a whole. It does not by itself
ensure the same coverage for every flare class or every type of observation.
Forecasting data are ordered in time, so the relation between calibration and
future samples also needs to be considered when interpreting the guarantee
\citep{xu2023spci}.

Class-conditional methods form calibration thresholds separately for each
class. This can target coverage within each class, but the estimate for a
class depends on the number of its calibration examples
\citep{vovk2005algorithmic,vovk2012conditional,ding2023classconditional}.
Strong flares are less common than weak or no-flare cases in many forecasting
datasets. A small number of strong-flare examples can therefore make
class-specific estimates less stable. This is one reason to report the
number of examples used for each class when evaluating uncertainty methods.

Flare classes have a natural order. Predicting an M-class event as C-class is
closer in severity than predicting it as an X-class event. Methods that treat
classes as unrelated labels may return a set with gaps, such as C and X but
not M. Ordinal conformal methods use the ordering to form contiguous sets or
intervals of classes \citep{lu2022ordinalaps,dey2023copoc}. Project 3 studies uncertainty sets for
ordered, imbalanced classes. Its method and evaluation are described in
Chapter~\ref{chap:project3}.

% ============================================================
\chapter{Solar Flare Target Formulation and Predictability}
\label{chap:project1}

\section{Motivation}

Machine-learning studies of solar-flare forecasting have commonly converted
future flare activity into a binary label by applying a fixed threshold to
GOES peak flux. This formulation supports a clear event/no-event decision, but
it discards how far a sample lies above or below the threshold. Consequently,
two windows with similar physical flux can receive different labels when they
fall on opposite sides of a cutoff, while windows with very different fluxes
on the same side receive the same label. This is a property of thresholded
label construction; it does not, by itself, establish where a classifier's
errors occur.

Prior work from this project has examined alternative flare metadata and
labeling schemes based on relative X-ray flux increases and cumulative
indices, and has evaluated these labels in classification settings
\citep{hong2023innovative,hong2023beyond,hong2023enhancing}. The work also
documents limitations of relying only on the maximum GOES class, including
the omission of background conditions and cumulative activity, and discusses
how background-flux variation can contribute to false alarms
\citep{hong2023beyond,hong2023enhancing}. The accompanying data release makes
the earlier label collection available for reuse \citep{hong2023data}.

This chapter extends that line of research from thresholded classification to
continuous regression. It formulates ten future XRS-B outcomes directly from
quality-filtered one-minute measurements and compares their predictability
from multi-channel solar observations. The target set includes maximum flux,
fluence, and maximum or total deviations from several pre-forecast background
levels. The central question is whether the choice of a continuous target
changes the predictive performance and uncertainty of otherwise comparable
models. The analysis does not assume that regression eliminates forecast
errors or that one target is universally preferable; it evaluates the
trade-off between target learnability and the physical quantity represented.

\section{Dataset and Preprocessing}

\subsection{Solar Observations}

The input observations are multi-channel images from NASA's Solar Dynamics
Observatory (SDO), using Atmospheric Imaging Assembly (AIA) and Helioseismic
and Magnetic Imager (HMI) channels. Each input image contains 13 channels
resampled to $224\times224$ pixels in a fixed order. The image
timestamps are matched to the forecast-start timestamps in the target tables.
The SDO observations and the published SuryaBench flare-forecasting benchmark
provide the image and sample-index context for this work
\citep{roy2026suryabench}.

\subsection{GOES XRS Measurements and Quality Control}

The ten primary regression targets are computed from one-minute GOES XRS-B
measurements rather than from individual entries in the NOAA flare catalog.
For each minute, the target-generation procedure uses a finite, positive
measurement that passes the instrument quality flag, preferring the scheduled
primary satellite and falling back to the scheduled secondary satellite when
needed \cite{TODO}. Missing minutes are not interpolated. A prediction window
uses all available valid samples, including when coverage is incomplete; a
target is missing only when its required prediction or background interval
contains no valid samples.
Event-catalog labels are outside the scope of the ten primary regression
outcomes; any event-based skill analysis will be reported separately.

\subsection{Data Splits}

All ten target conditions use the same training, validation, and test
partitions and the same validation-based model-selection procedure. Before
filtering for target validity and image availability, these partitions contain
74,760 training rows, 3,672 validation rows, and 43,848 test rows; the usable
sample count can differ by target. The final manuscript will report the UTC
date ranges and usable counts from the corresponding data records. Any overlap
between 24-hour target windows at the temporal boundary will be reported from
an audit of these same partitions.

\section{Target Definitions}

Let $s$ denote the forecast start and let $F_i(s)$ be the $i$th eligible
one-minute XRS-B flux sample. The half-open forecast window is
$[s,s+24\,\mathrm{h})$. It contains $N_p(s)$ valid samples, where
$N_p(s)$ may be less than the 1,440 expected minute bins. The ten outcomes
comprise two background-independent targets and eight signed background
deviation targets.

\subsection{Background-Independent Targets}

The maximum-flux target is
\begin{equation}
 y_{\mathrm{max\text{-}xflux}}(s)
 = \max_{i\in\mathcal{P}(s)} F_i(s),
\end{equation}
where $\mathcal{P}(s)$ indexes the valid samples in the forecast window.
The fluence target is the unscaled one-minute sample sum
\begin{equation}
 y_{\mathrm{fluence}}(s)
 = \sum_{i\in\mathcal{P}(s)} F_i(s).
\end{equation}
The sum is not multiplied by 60 seconds and is therefore reported as an
unscaled sample sum rather than as a time-integrated SI fluence.

\subsection{Background-Deviation Targets}

For background-window length $W\in\{2,24\}$ hours, define the valid samples in
the immediately preceding half-open interval $[s-W,s)$ as
$R_1,\ldots,R_{N_b}$. For lower-tail percentage $q\in\{10,50\}$, let
$k=\max(1,\lceil qN_b/100\rceil)$ and define the background level as the mean
of the $k$ smallest valid samples when $N_b>0$:
\begin{equation}
 B_{W,q}(s)=\frac{1}{k}\sum_{j=1}^{k} R_{(j)},
 \qquad R_{(1)}\le\cdots\le R_{(N_b)}.
\end{equation}
This background level is fixed over the forecast window. For each of the four
$(W,q)$ combinations, define the signed sample deviations
$D_i^{W,q}(s)=F_i(s)-B_{W,q}(s)$. The two corresponding targets are
\begin{align}
 y_{\mathrm{max\text{-}dev},W,q}(s)
 &=\max_{i\in\mathcal{P}(s)}D_i^{W,q}(s),\\
 y_{\mathrm{total\text{-}dev},W,q}(s)
 &=\sum_{i\in\mathcal{P}(s)}D_i^{W,q}(s).
\end{align}
Together, $(W,q)\in\{(2,10),(2,50),(24,10),(24,50)\}$ yields eight
background-deviation targets. The total deviation is an unscaled one-minute
sum; negative deviations are retained and are not clipped to zero. If the
forecast window has no valid samples, all ten targets are missing. If a
background interval has no valid samples, only the two deviation targets
depending on that background are missing.

\begin{table}[ht]
    \centering
    \caption{The ten continuous regression targets used in this chapter.}
    \label{tab:project1-target-definitions}
    \small
    \begin{tabular}{p{0.25\linewidth}p{0.67\linewidth}}
        \toprule
        Target family & Definition \\
        \midrule
        Maximum flux & Maximum XRS-B flux during the 24-hour forecast window. \\
        Fluence & Sum of valid one-minute XRS-B flux samples during the forecast window. \\
        Maximum deviation & Maximum forecast-window flux minus the preceding-window background level; four settings combine 2- or 24-hour background windows with 10\% or 50\% lower-tail means. \\
        Total deviation & Sum of signed forecast-window flux deviations from the same four background levels. \\
        \bottomrule
    \end{tabular}
    \normalsize
\end{table}

\section{Model and Training Protocol}

\subsection{Multi-Channel Image Model}

The model uses eight AIA channels and five HMI channels, each resampled to
$224\times224$ pixels. The target comparisons use the same compact
convolutional neural network architecture. Input normalization is estimated
using training data only. Candidate hyperparameter settings are compared
under the common validation procedure; the selected settings and number of
independent training runs will be reported with the results.

\subsection{Quantile Regression}

The model predicts the $0.05$, $0.50$, and $0.95$ conditional quantiles using
the pinball loss. Target values are transformed using an inverse-hyperbolic-
sine transformation followed by standardization; both transformations are
estimated from training data only. Predictions are converted back to the
target's original scale before evaluation. Model selection uses validation
data, and test labels are reserved for final evaluation.

\subsection{Evaluation Metrics}

The primary point-prediction metrics are mean absolute error (MAE), mean
squared error (MSE), root mean squared error (RMSE), Spearman rank correlation,
and coefficient of determination ($R^2$), calculated from the predicted
median ($q_{0.50}$). The pinball loss remains the model's training objective,
but is not a reported comparison metric. Binary skill metrics will not be used
as primary regression metrics; any threshold-based event analysis will be
reported separately with its threshold and event denominator.

\section{Preliminary Results}

The final comparison will use the common split and validation-based
model-selection protocol described above. The table below is a reporting
template; values will be populated from saved test-prediction artifacts only
after test evaluation artifacts are available and their run provenance is
checked.

\begin{table}[ht]
    \centering
    \caption{Planned comparison of regression performance across the ten target conditions.}
    \label{tab:target-comparison}
    \footnotesize
    \begin{tabular}{>{\raggedright\arraybackslash}p{0.34\linewidth}ccccc}
        \toprule
        Target condition & MAE & MSE & RMSE & Spearman $\rho$ & $R^2$ \\
        \midrule
        Maximum XRS-B flux & -- & -- & -- & -- & -- \\
        Fluence & -- & -- & -- & -- & -- \\
        Max. deviation (2 h, 10\%) & -- & -- & -- & -- & -- \\
        Total deviation (2 h, 10\%) & -- & -- & -- & -- & -- \\
        Max. deviation (2 h, 50\%) & -- & -- & -- & -- & -- \\
        Total deviation (2 h, 50\%) & -- & -- & -- & -- & -- \\
        Max. deviation (24 h, 10\%) & -- & -- & -- & -- & -- \\
        Total deviation (24 h, 10\%) & -- & -- & -- & -- & -- \\
        Max. deviation (24 h, 50\%) & -- & -- & -- & -- & -- \\
        Total deviation (24 h, 50\%) & -- & -- & -- & -- & -- \\
        \bottomrule
    \end{tabular}
    \normalsize
    \par\smallskip
    \footnotesize For deviation targets, the first value is the background-window duration and the second is the lower-tail percentage used to estimate the background level.
\end{table}

\section{Discussion}

\section{Remaining Work}

% ============================================================
\chapter{Multi-channel Solar Flare Forecasting}
\label{chap:project2}

\section{Motivation}

\section{Research Questions}

\section{Input Modalities}

\subsection{HMI}
\subsection{AIA}
\subsection{Single-channel and Multi-channel Configurations}

\section{Experimental Design}

\subsection{Prediction Windows}

% e.g. 8 h, 12 h, 24 h

\subsection{Image Resolution}

% e.g. 224 x 224 and 512 x 512

\subsection{Model Capacity}

% Describe the model-size variants included in Project 2's experiments.

\subsection{Target Variants}

\section{Preliminary Results}

\section{Planned Experiments}

\section{Expected Findings}

% ============================================================
\chapter{Uncertainty Quantification for Ordinal Prediction}
\label{chap:project3}

\section{Motivation}

Ordinal classifiers predict among categories with a meaningful order, such as
increasing disease severity, age groups, or solar flare classes. A point
prediction hides both which alternative classes remain plausible and how far
those alternatives lie from the predicted class. Prediction sets expose this
uncertainty, but a useful set for an ordinal task should also respect the
ordering: a set containing two separated classes while omitting the classes
between them is difficult to interpret as a range of possible outcomes.

Conformal prediction can provide finite-sample coverage guarantees under
exchangeability, but standard classification methods target marginal
coverage and treat labels as nominal. Under class imbalance, their aggregate
coverage can conceal poor coverage of rare classes. Ordinal conformal methods
such as OAPS and COPOC account for label order, but their standard guarantees
are marginal. This project builds on the published Ordinal Conformalized
Quantile Regression (OCQR) framework, which combines class-specific
calibration with contiguous prediction sets. The method is designed to
address both class-wise reliability and ordinal interpretability
\citep{hong2026ocqr}.

\section{Problem Formulation}

Let \((X,Z,Y)\) denote an example, where \(X\in\mathcal{X}\) is the model
input, \(Z\in\mathbb{R}\) is a numeric target used for quantile regression,
and \(Y\in\mathcal{Y}=\{0,\ldots,K-1\}\) is the ordinal label. Fixed
boundaries \(b_1<\cdots<b_{K-1}\) define ordered bins
\(B_0,\ldots,B_{K-1}\); the first and last bins extend to negative and
positive infinity. The numeric representation and bin boundaries are fixed
before calibration. Specifically, \(B_0=(-\infty,b_1)\),
\(B_k=[b_k,b_{k+1})\) for \(1\leq k<K-1\), and
\(B_{K-1}=[b_{K-1},+\infty)\). The coverage argument assumes
target-label consistency:
\begin{equation}
Y=k \quad\Longrightarrow\quad Z\in B_k,
\label{eq:ocqr-target-label-consistency}
\end{equation}
for every \(k\in\mathcal{Y}\). The goal is to construct a set
\(C(X)\subseteq\mathcal{Y}\) that contains the true ordinal class at the
desired rate within each class and is nonempty and contiguous. A set is
contiguous if it contains every class between its smallest and largest
members. The numeric target is used to fit quantiles and compute calibration
scores; the observed class label defines the ordinal prediction task and the
true-label calibration group \citep{hong2026ocqr}.

\section{Conformal Prediction Background}

\subsection{Split Conformal Prediction}

Split conformal prediction separates model fitting from uncertainty
calibration. After fitting the predictive model and fixing other
data-dependent choices on training and validation data, a held-out calibration
set is used to calculate nonconformity scores. For scores
\(S_1,\ldots,S_n\) and target miscoverage \(\alpha\), the split-conformal
threshold is the order statistic
\begin{equation}
\widehat q_{1-\alpha}
=S_{(\lceil(n+1)(1-\alpha)\rceil)},
\label{eq:ocqr-split-threshold}
\end{equation}
where the calibration scores are augmented with one \(+\infty\) value so the
requested rank is defined even for small calibration samples. A test label is
included when its score does not exceed the calibrated threshold. If the
calibration observations and a future test observation are exchangeable, the
resulting set has marginal coverage at least \(1-\alpha\). This guarantee
averages over the population and does not require a parametric model for the
data distribution \citep{vovk1999cp,romano2020classification}.

For classification, a score is computed for each candidate class, and the
prediction set contains the classes whose scores pass the threshold. LAC and
APS are examples of such methods. Their usual forms treat the labels as
nominal, so a valid set may include separated classes while omitting an
intermediate class \citep{sadinle2019least,angelopoulos2020uncertainty,
romano2020classification}.

\subsection{Class-conditional Coverage}

Marginal coverage does not ensure that each class has the same coverage. A
class-conditional procedure instead targets
\(\Pr\{Y\in C(X)\mid Y=k\}\geq 1-\alpha\) for each class \(k\). Mondrian
conformal prediction obtains group-specific guarantees by splitting the
calibration observations into groups and estimating a separate threshold
within each group. When the groups are defined by true labels, this is
class-conditional calibration. Under exchangeability within each class, the
class-specific threshold gives a finite-sample coverage guarantee for a
future observation from that class \citep{vovk2005algorithmic,
vovk2012conditional}.

\subsection{Challenges Under Class Imbalance}

Class-specific calibration uses only the calibration examples belonging to
the class being evaluated. Rare classes therefore have fewer scores from
which to estimate a quantile. The finite-sample rank can exceed the number of
finite class-specific scores; in that case the augmented order statistic is
\(+\infty\), and the class is included for every input. This conservative
behavior preserves the coverage construction but can enlarge prediction
sets. It also makes class counts and per-class coverage necessary parts of
the evaluation, alongside marginal coverage and set size
\citep{ding2023classconditional}.

\section{Proposed Method}

OCQR combines quantile regression with true-label Mondrian calibration. At
inference, it evaluates every candidate class using that class's own
conformal correction, maps the resulting numeric acceptance region to the
candidate's ordinal bin, and then closes the accepted labels into a
nonempty contiguous interval. The method does not require an
ordinal-specific output layer in the predictive backbone \citep{hong2026ocqr}.

\subsection{Quantile Regression}

For target miscoverage \(\alpha\), OCQR estimates lower and upper conditional
quantiles of \(Z\), at levels \(\tau_L=\alpha/2\) and
\(\tau_U=1-\alpha/2\). Quantile regression is fitted with the pinball loss
\begin{equation}
\rho_{\tau}(r)=\max\{\tau r,(\tau-1)r\},
\label{eq:ocqr-pinball-loss}
\end{equation}
where \(r=z-\widehat q_{\tau}(x)\). To handle quantile crossing, the
predicted endpoints are ordered as
\begin{equation}
\widehat\ell(x)=\min\{\widehat q_{\tau_L}(x),\widehat q_{\tau_U}(x)\},
\qquad
\widehat u(x)=\max\{\widehat q_{\tau_L}(x),\widehat q_{\tau_U}(x)\}.
\label{eq:ocqr-ordered-endpoints}
\end{equation}
The conformalized-quantile-regression score is
\begin{equation}
s(x,z)=\max\{\widehat\ell(x)-z,\;z-\widehat u(x)\}.
\label{eq:ocqr-score}
\end{equation}
For a correction \(q\), its numeric acceptance region is
\(A_q(x)=\{z:s(x,z)\leq q\}\). For finite \(q\), this is
\([\widehat\ell(x)-q,\widehat u(x)+q]\) when the lower endpoint does not
exceed the upper endpoint, and is empty otherwise. For \(q=+\infty\), it
is defined as \(\mathbb{R}\) \citep{romano2019cqr,hong2026ocqr}.

\subsection{Calibration}

Let \(\mathcal{D}_{\mathrm{cal}}=\{(X_i,Z_i,Y_i)\}_{i=1}^n\) be the
calibration set. For class \(k\), define
\(\mathcal{I}_k=\{i:Y_i=k\}\) and \(N_k=|\mathcal{I}_k|\). The calibration
scores for that class are \(S_{k,i}=s(X_i,Z_i)\), for
\(i\in\mathcal{I}_k\). OCQR computes the finite-sample rank and class-specific
correction as
\begin{equation}
r_k=\left\lceil(N_k+1)(1-\alpha)\right\rceil,
\qquad
q_k=\widetilde S_{k,(r_k)},
\label{eq:ocqr-class-correction}
\end{equation}
where \(\widetilde S_{k,(r_k)}\) is the \(r_k\)-th order statistic of the
class-\(k\) scores augmented with \(+\infty\). The model, target
representation, bin boundaries, and method choices are fixed before this
calibration step; calibration data are used only to estimate the corrections.
When the class has too few calibration observations to support the requested
finite rank, \(q_k=+\infty\), so the corresponding class is always retained.
This rule prioritizes the finite-sample coverage construction over set-size
efficiency for classes with very limited calibration support
\citep{hong2026ocqr}.

\subsection{Ordinal Contiguity}

At test time, the true class is unknown, so OCQR tests each candidate \(k\)
with its own correction \(q_k\). Its raw set is
\begin{equation}
C_{\mathrm{raw}}(x)
=\{k\in\mathcal{Y}:A_{q_k}(x)\cap B_k\neq\varnothing\}.
\label{eq:ocqr-raw-set}
\end{equation}
Thus a candidate is accepted when its calibrated numeric region intersects
the numeric bin associated with that class. The resulting set may have gaps.
OCQR first replaces an empty raw set by the full label set, then fills every
class between the smallest and largest accepted labels:
\begin{equation}
C(x)=\{j\in\mathcal{Y}:\min\widetilde C_{\mathrm{raw}}(x)
\leq j\leq\max\widetilde C_{\mathrm{raw}}(x)\},
\label{eq:ocqr-final-set}
\end{equation}
where \(\widetilde C_{\mathrm{raw}}(x)=C_{\mathrm{raw}}(x)\) when the raw
set is nonempty and \(\mathcal{Y}\) otherwise. The fallback and ordinal hull
only add labels, so the final set is nonempty and contiguous and contains the
raw set \citep{hong2026ocqr}.

\subsection{Rare-class Handling}

OCQR does not pool or borrow calibration scores across distinct classes in
its canonical class-conditional procedure. Each class receives its own
correction, and the augmented order statistic handles insufficient support
conservatively. A pooled correction is included as an ablation, but does not
provide the same class-specific calibration argument. Accordingly, the
evaluation reports the calibration and test counts for every class, the
per-class coverage estimates, and the size of the resulting sets. When a
class-specific correction is infinite, its guaranteed inclusion can make
sets less informative; this is a limitation of sparse calibration support,
not a reason to interpret the class as well resolved
\citep{ding2023classconditional,hong2026ocqr}.

\section{Theoretical Properties}

Let \(\mathcal{F}\) contain all information used to fit or select the
procedure before calibration, including the fitted model, target
representation, bin boundaries, hyperparameters, Mondrian groups, and
post-processing rules. Assume these components are fixed without using
calibration or test observations. For each class \(k\), assume that the
class-\(k\) calibration scores and the corresponding test score are
exchangeable conditional on \(\mathcal{F}\), and assume
\cref{eq:ocqr-target-label-consistency}.

\noindent\textbf{Theorem (Class-conditional coverage).}
Under these assumptions, the OCQR prediction set satisfies, for every
\(k\in\mathcal{Y}\),
\begin{equation}
\Pr\{Y_{n+1}\in C(X_{n+1})\mid\mathcal{F},Y_{n+1}=k\}
\geq 1-\alpha.
\label{eq:ocqr-conditional-coverage}
\end{equation}
Consequently,
\begin{equation}
\Pr\{Y_{n+1}\in C(X_{n+1})\mid Y_{n+1}=k\}
\geq 1-\alpha.
\label{eq:ocqr-marginalized-conditional-coverage}
\end{equation}
The final set is nonempty and contiguous.

\noindent\textit{Argument.}
Conditional on \(\mathcal{F}\), the class count, and test class \(k\), the
augmented split-conformal rank gives
\(s(X_{n+1},Z_{n+1})\leq q_k\) with probability at least
\(1-\alpha\). Therefore \(Z_{n+1}\in A_{q_k}(X_{n+1})\). By
target-label consistency, \(Z_{n+1}\in B_k\), so the two sets intersect and
\(k\in C_{\mathrm{raw}}(X_{n+1})\). The fallback and ordinal hull only add
labels, preserving this inclusion and proving the stated coverage. Averaging
over the class-specific calibration count and over \(\mathcal{F}\) yields
the two displayed guarantees. The guarantee depends on within-class
exchangeability and target-label consistency; it does not assert conditional
coverage for every input \(X=x\) \citep{hong2026ocqr}.

\section{Experimental Evaluation}

The published evaluation compares OCQR with nominal and ordinal conformal
baselines on three five-class tasks. It uses fixed train, validation,
calibration, and test partitions, keeps those partitions unchanged across
five training seeds, and sets the nominal coverage to \(1-\alpha=0.9\).
The evaluation reports coverage and set efficiency so rare-class reliability
is not obscured by a single population average \citep{hong2026ocqr}.

\subsection{Datasets}

The evaluation uses UTKFace for age estimation, RetinaMNIST for
diabetic-retinopathy severity grading, and SuryaBench for solar flare
forecasting \citep{utkface-zhifei2017cvpr,medmnistv2,roy2026suryabench}.
UTKFace uses chronological age in years as its numeric target and five fixed
age bins: below 20, 20--39, 40--59, 60--79, and 80 years or older. Its
23,708 images are split into 14,224 training, 2,371 validation, 4,742
calibration, and 2,371 test examples. The split is at the image level, not
the subject level.

RetinaMNIST provides five ordered severity labels but no continuous
measurement. It therefore uses the fixed class-index embedding \(Z=Y\), with
bin boundaries at 0.5, 1.5, 2.5, and 3.5. Its official validation and test
partitions contain 120 and 400 images; the official 1,080-image training
partition is further divided into 756 training and 324 calibration examples.

The solar flare task uses a 13-channel SDO image stack to predict maximum
GOES X-ray flux over the following 24 hours. Its numeric target is
\(Z=\log_{10}(\texttt{max\_intensity})+9\), with ordered classes FQ/A, B, C,
M, and X and transformed bin boundaries 2, 3, 4, and 5. Within each year
from 2010 to 2019, observations from January 1--14 and February 1--14 are
used for calibration, January 15--31 for validation, and February 15--
December 31 for training. The future test set contains retained observations
from 2020--2024. This chronology makes the solar test a temporal-extrapolation
evaluation; it does not establish exchangeability between calibration and
future test observations \citep{roy2026suryabench,hong2026ocqr}.

\subsection{Baselines}

OCQR is compared with the Least Ambiguous set-valued Classifier (LAC),
Adaptive Prediction Sets (APS), Ordinal Adaptive Prediction Sets (OAPS),
and Conformal Prediction Sets for Ordinal Classification (COPOC)
\citep{sadinle2019least,angelopoulos2020uncertainty,
romano2020classification,lu2022ordinalaps,dey2023copoc}. LAC, APS, and OAPS
use ResNet-18 Softmax classifiers; COPOC uses a separately trained
ResNet-18 with its ordinal-specific unimodal output. Five OCQR ablations
isolate the contribution of pooled rather than class-specific calibration,
the ordinal hull, the empty-set fallback, the combination of hull and
fallback, and replacing negative conformity scores by zero
\citep{hong2026ocqr}.

\subsection{Evaluation Metrics}

For class \(k\), empirical class-conditional coverage is
\begin{equation}
\widehat{\mathrm{Cov}}_k=
\frac{1}{n_k}\sum_{i:Y_i=k}\mathbf{1}\{Y_i\in C(X_i)\},
\label{eq:ocqr-empirical-class-coverage}
\end{equation}
where \(n_k\) is the number of test examples in that class. The evaluation
reports marginal coverage, macro coverage (the mean of the class-specific
estimates), worst-class coverage (their minimum), mean set size, and the
fraction of examples receiving the full label set. It also reports empty-set
and fragmented-set rates for the ablations. Results are means and standard
deviations across five training seeds; the fixed splits mean these standard
deviations describe training-seed variation, not variability across
alternative partitions \citep{hong2026ocqr}.

\subsection{Ablation Studies}

The pooled-calibration ablation tests whether class-specific corrections
account for rare-class performance. The no-hull and raw-set ablations test
whether ordinal closure changes contiguity, while the no-fallback and raw-set
ablations test whether the fallback prevents empty predictions. The
nonnegative-score variant tests sensitivity to the signed CQR score. These
ablations are evaluated on all three datasets using the same primary metrics
and additionally record empty and fragmented-set rates \citep{hong2026ocqr}.

\section{Preliminary Results}

Across the three datasets, OCQR occupied an intermediate point in the
coverage--set-size tradeoff: it improved worst-class coverage over the more
compact LAC and OAPS baselines while producing smaller sets than APS and
COPOC in the reported comparisons. On UTKFace, OCQR achieved marginal
coverage of \(98.7\%\pm0.2\%\), macro class coverage of
\(98.1\%\pm0.4\%\), and worst-class coverage of \(96.7\%\pm0.8\%\), with a
mean set size of \(2.35\pm0.06\). LAC and OAPS had mean set sizes near 1.6
classes, but their worst-class coverage was 66.4\% and 68.7\%, respectively.

On RetinaMNIST, OCQR achieved \(95.8\%\pm1.2\%\) marginal,
\(93.3\%\pm1.3\%\) macro, and \(81.4\%\pm4.7\%\) worst-class coverage,
with mean set size \(4.04\pm0.18\). APS and COPOC achieved higher
worst-class coverage (98.7\% and 95.0\%) but had mean set sizes of 4.91 and
4.74 out of five classes. OCQR's class-4 coverage was \(82.0\%\pm5.7\%\),
below the nominal target; this empirical value on a fixed test set does not
contradict a probabilistic finite-sample guarantee.

On the chronological solar test period, OCQR achieved
\(89.6\%\pm5.4\%\) marginal, \(88.4\%\pm6.2\%\) macro, and
\(73.0\%\pm17.8\%\) worst-class coverage, with mean set size
\(2.60\pm0.11\). Its X-class coverage was \(86.1\%\pm6.6\%\), while
LAC and OAPS attained 10.0\% and 5.8\%, respectively. OCQR's weakest solar
class was B, at \(73.2\%\pm18.2\%\). These results are empirical under
temporal extrapolation and are not evidence that the exchangeability
assumption holds between the calibration and test periods.

In the pooled-calibration ablation, worst-class coverage decreased from
81.4\% to 53.0\% on RetinaMNIST, from 96.7\% to 87.2\% on UTKFace, and
from 73.0\% to 43.9\% on solar forecasting. Removing the ordinal hull
produced fragmented sets at \(0.7\%\pm1.0\%\) on RetinaMNIST, whereas the
canonical hull eliminated those cases. No empty sets occurred in the reported
experiments, so removing the fallback did not change the observed metrics.
The results support the role of class-specific calibration in worst-class
coverage and of the hull and fallback as structural safeguards
\citep{hong2026ocqr}.

\section{Remaining Theoretical and Experimental Work}

% TODO (author input): Specify which dissertation contributions remain beyond
% the published OCQR method and experiments. The source paper identifies
% additional nominal coverage levels, alternative target representations,
% more distribution-shift settings, and additional ordinal tasks as future
% directions, but it does not establish which of these are part of this
% proposal's committed work plan.
% TODO (author input): State whether the dissertation will study guarantees
% for temporally dependent or shifted solar data. The reported chronological
% solar results are empirical temporal-extrapolation results and do not
% verify the exchangeability assumption behind the finite-sample guarantee.

% ============================================================
\chapter{Integrated Discussion}
\label{chap:discussion}

\section{Connections Across the Three Projects}

% Explain the common theme:
% reliable scientific prediction under challenging labels,
% imbalance, and uncertainty.

\section{Scientific and Methodological Contributions}

\section{Limitations}

% ============================================================
\chapter{Proposed Work and Timeline}
\label{chap:timeline}

\section{Remaining Work}

\begin{enumerate}
    \item Complete additional forecasting-window experiments.
    \item Complete multi-channel and higher-resolution experiments.
    \item Finalize the ordinal uncertainty quantification method.
    \item Complete theoretical analysis.
    \item Run full benchmark and ablation experiments.
    \item Prepare dissertation manuscripts and final dissertation.
\end{enumerate}

\section{Timeline}

\begin{table}[ht]
    \centering
    \caption{Tentative dissertation timeline.}
    \label{tab:timeline}
    \begin{tabular}{p{0.25\textwidth}p{0.65\textwidth}}
        \toprule
        Period & Planned Work \\
        \midrule
        Fall 2026 &
        Dissertation proposal; complete core solar-flare experiments;
        finalize Project 3 methodology. \\

        Spring 2027 &
        Complete Project 2 and Project 3 experiments; prepare manuscripts;
        integrate dissertation chapters. \\

        Summer 2027 &
        Complete dissertation writing, committee review, and dissertation
        defense preparation. \\
        \bottomrule
    \end{tabular}
\end{table}

\section{Risks and Contingency Plans}

% Optional but useful in a proposal:
% - compute constraints
% - data limitations
% - theoretical assumptions
% - fallback experimental scope

% ============================================================
\chapter{Conclusion}
\label{chap:conclusion}

% Summarize:
% - dissertation theme
% - completed work
% - remaining work
% - expected contribution

% ============================================================
% REFERENCES
% ============================================================
\bibliographystyle{unsrtnat}
\bibliography{references}

\end{document}
